\begin{thebibliography}{99}
\bibitem{AFM15} N. Anantharaman, C. Fermanian Kammerer \& F. Macià - “Semiclassical completely integrable systems: long-time dynamics and observability via two-microlocal Wigner measures” , Amer. J. Math. 137 (2015) no. 3, p. 577-638
\bibitem{AM14} N. Anantharaman \& F. Macià - “Semiclassical measures for the Schrödinger equation on the torus” , J. Eur. Math. Soc. (JEMS) 16 (2014) no. 6, p. 1253-1288
\bibitem{Ba} H. Bahouri - “Non prolongement unique des solutions d’opérateurs ‘somme de carrés’” , Ann. Inst. Fourier (Grenoble) 36 (1986) no. 4, p. 137-155
\bibitem{BFG2} H. Bahouri, C. Fermanian Kammerer \& I. Gallagher - “Dispersive estimates for the Schrödinger operator on step-2 stratified Lie groups” , Anal. PDE 9 (2016) no. 3, p. 545-574
\bibitem{BGX} H. Bahouri, P. Gérard \& C.-J. Xu - “Espaces de Besov et estimations de Strichartz généralisées sur le groupe de Heisenberg” , J. Anal. Math. 82 (2000), p. 93-118
\bibitem{BC} K. Beauchard \& P. Cannarsa - “Heat equation on the Heisenberg group: observability and applications” , J. Differential Equations 262 (2017) no. 8, p. 4475-4521
\bibitem{BDE} K. Beauchard, J. Dardé \& S. Ervedoza - “Minimal time issues for the observability of Grushin-type equations” , Ann. Inst. Fourier (Grenoble) 70 (2020) no. 1, p. 247-312
\bibitem{BLU} A. Bonfiglioli, E. Lanconelli \& F. Uguzzoni - Stratified Lie groups and potential theory for their sub-Laplacians , Springer Monographs in Math. , Springer, Berlin, 2007
\bibitem{Bon69} J.-M. Bony - “Principe du maximum, inégalité de Harnack et unicité du problème de Cauchy pour les opérateurs elliptiques dégénérés” , Ann. Inst. Fourier (Grenoble) 19 (1969) no. 1, p. 277-304
\bibitem{BS} N. Burq \& C. Sun - “Time optimal observability for Grushin Schrödinger equation” , 2019, to appear in Anal. PDE
\bibitem{BZ} N. Burq \& M. Zworski - “Control for Schrödinger operators on tori” , Math. Res. Lett. 19 (2012) no. 2, p. 309-324
\bibitem{CFM19} V. Chabu, C. Fermanian Kammerer \& F. Macià - “Semiclassical analysis of dispersion phenomena” , in Analysis and partial differential equations: perspectives from developing countries , Springer Proc. Math. Stat. , vol. 275 , Springer, Cham, 2019, p. 84-108
\bibitem{CFM20} V. Chabu, C. Fermanian Kammerer \& F. Macià - “Wigner measures and effective mass theorems” , Ann. H. Lebesgue 3 (2020), p. 1049-1089
\bibitem{CR} M. Combescure \& D. Robert - Coherent states and applications in mathematical physics , Theoretical and Math. Physics , Springer, Dordrecht, 2012
\bibitem{CG} L. J. Corwin \& F. P. Greenleaf - Representations of nilpotent Lie groups and their applications. Part I. Basic theory and examples , Cambridge Studies in Advanced Math. , vol. 18 , Cambridge University Press, Cambridge, 1990
\bibitem{DK} M. Duprez \& A. Koenig - “Control of the Grushin equation: non-rectangular control region and minimal time” , ESAIM Control Optim. Calc. Var. 26 (2020), article ID 3, 18 pages
\bibitem{Fermanian_note1} C. Fermanian Kammerer - “Mesures semi-classiques 2-microlocales” , C. R. Acad. Sci. Paris Sér. I Math. 331 (2000) no. 7, p. 515-518
\bibitem{Fermanian_Note2} C. Fermanian Kammerer - “Analyse à deux échelles d’une suite bornée de $L^2$ sur une sous-variété du cotangent” , Comptes Rendus Mathématique 340 (2005) no. 4, p. 269-274
\bibitem{FF} C. Fermanian Kammerer \& V. Fischer - “Quantum evolution and sub-Laplacian operators on groups of Heisenberg type” , 2019, to appear in J. Spectral Theory
\bibitem{FF1} C. Fermanian Kammerer \& V. Fischer - “Semi-classical analysis on H-type groups” , Sci. China Math. 62 (2019) no. 6, p. 1057-1086
\bibitem{FF18} C. Fermanian Kammerer \& V. Fischer - “Defect measures on graded Lie groups” , Ann. Scuola Norm. Sup. Pisa Cl. Sci. (5) 21 (2020), p. 207-291
\bibitem{FG02} C. Fermanian Kammerer \& P. Gérard - “Mesures semi-classiques et croisement de modes” , Bull. Soc. math. France 130 (2002) no. 1, p. 123-168
\bibitem{FL08} C. Fermanian Kammerer \& C. Lasser - “Propagation through generic level crossings: a surface hopping semigroup” , SIAM J. Math. Anal. 40 (2008) no. 1, p. 103-133
\bibitem{FR} V. Fischer \& M. Ruzhansky - Quantization on nilpotent Lie groups , Progress in Math. , vol. 314 , Birkhäuser/Springer, Cham, 2016
\bibitem{gerard_X} P. Gérard - “Mesures semi-classiques et ondes de Bloch” , in Séminaire sur les Équations aux Dérivées Partielles, 1990–1991 , École Polytechnique, Palaiseau, 1991, Exp. No. XVI, 19 p.
\bibitem{gerard_91} P. Gérard - “Microlocal defect measures” , Comm. Partial Differential Equations 16 (1991) no. 11, p. 1761-1794
\bibitem{gerardleichtnam} P. Gérard \& É. Leichtnam - “Ergodic properties of eigenfunctions for the Dirichlet problem” , Duke Math. J. 71 (1993) no. 2, p. 559-607
\bibitem{GMMP} P. Gérard, P. A. Markowich, N. J. Mauser \& F. Poupaud - “Homogenization limits and Wigner transforms” , Comm. Pure Appl. Math. 50 (1997) no. 4, p. 323-379, Erratum: Ibid. 53 (2000), no. 2, p. 280–281
\bibitem{hag} G. A. Hagedorn - “Semiclassical quantum mechanics. I. The $\hbar \rightarrow 0$ limit for coherent states” , Comm. Math. Phys. 71 (1980) no. 1, p. 77-93
\bibitem{HMR} B. Helffer, A. Martinez \& D. Robert - “Ergodicité et limite semi-classique” , Comm. Math. Phys. 109 (1987) no. 2, p. 313-326
\bibitem{Hor67} L. Hörmander - “Hypoelliptic second order differential equations” , Acta Math. 119 (1967), p. 147-171
\bibitem{J} S. Jaffard - “Contrôle interne exact des vibrations d’une plaque rectangulaire” , Portugal. Math. 47 (1990) no. 4, p. 423-429
\bibitem{Kap80} A. Kaplan - “Fundamental solutions for a class of hypoelliptic PDE generated by composition of quadratic forms” , Trans. Amer. Math. Soc. 258 (1980) no. 1, p. 147-153
\bibitem{Koe} A. Koenig - “Non-null-controllability of the Grushin operator in 2D” , Comptes Rendus Mathématique 355 (2017) no. 12, p. 1215-1235
\bibitem{LT05} C. Lasser \& S. Teufel - “Propagation through conical crossings: an asymptotic semigroup” , Comm. Pure Appl. Math. 58 (2005) no. 9, p. 1188-1230
\bibitem{Laur} C. Laurent \& M. Léautaud - “Tunneling estimates and approximate controllability for hypoelliptic equations” , 2017, to appear in Mem. Amer. Math. Soc.
\bibitem{Leb2} G. Lebeau - “Control for hyperbolic equations” , in Journées “Équations aux Dérivées Partielles” (Saint-Jean-de-Monts, 1992) , École Polytechnique, Palaiseau, 1992, p. 24
\bibitem{Leb} G. Lebeau - “Contrôle de l’équation de Schrödinger” , J. Math. Pures Appl. (9) 71 (1992) no. 3, p. 267-291
\bibitem{lebeau} G. Lebeau - “Équation des ondes amorties” , in Algebraic and geometric methods in mathematical physics (Kaciveli, 1993) , Math. Phys. Stud. , vol. 19 , Kluwer Acad. Publ., Dordrecht, 1996, p. 73-109
\bibitem{Let} C. Letrouit - “Subelliptic wave equations are never observable” , 2020
\bibitem{LS} C. Letrouit \& C. Sun - “Observability of Baouendi-Grushin-type equations through resolvent estimates” , 2020, to appear in J. Inst. Math. Jussieu
\bibitem{JLL} J.-L. Lions - Contrôlabilité exacte, perturbations et stabilisation de systèmes distribués. Tome 1. Contrôlabilité exacte , Recherches en Mathématiques Appliquées , vol. 8 , Masson, Paris, 1988
\bibitem{LP93} P.-L. Lions \& T. Paul - “Sur les mesures de Wigner” , Rev. Mat. Iberoamericana 9 (1993) no. 3, p. 553-618
\bibitem{MaciaTorus} F. Macià - “High-frequency propagation for the Schrödinger equation on the torus” , J. Funct. Anal. 258 (2010) no. 3, p. 933-955
\bibitem{Mac11} F. Macià - “The Schrödinger flow in a compact manifold: high-frequency dynamics and dispersion” , in Modern aspects of the theory of partial differential equations , Oper. Theory Adv. Appl. , vol. 216 , Birkhäuser/Springer Basel AG, Basel, 2011, p. 275-289
\bibitem{MaciaLille} F. Macià - “High-frequency dynamics for the Schrödinger equation, with applications to dispersion and observability” , in Nonlinear optical and atomic systems , Lect. Notes in Math. , vol. 2146 , Springer, Cham, 2015, p. 275-335
\bibitem{MacRiv18} F. Macià \& G. Rivière - “Two-microlocal regularity of quasimodes on the torus” , Anal. PDE 11 (2018) no. 8, p. 2111-2136
\bibitem{MacRiv19} F. Macià \& G. Rivière - “Observability and quantum limits for the Schrödinger equation on $\mathbb{S}^d$” , in Probabilistic methods in geometry, topology and spectral theory , American Mathematical Society, Providence, RI, 2019, p. 139-153
\bibitem{MillerThesis} L. Miller - Propagation d’ondes semi-classiques à travers une interface et mesures 2-microlocales , Ph. D. Thesis, École Polytechnique, Palaiseau, 1996
\bibitem{NierScat} F. Nier - “A semi-classical picture of quantum scattering” , Ann. Sci. École Norm. Sup. (4) 29 (1996) no. 2, p. 149-183
\bibitem{pedersen} N. V. Pedersen - “Matrix coefficients and a Weyl correspondence for nilpotent Lie groups” , Invent. Math. 118 (1994) no. 1, p. 1-36
\bibitem{tartar} L. Tartar - “$H$-measures, a new approach for studying homogenisation, oscillations and concentration effects in partial differential equations” , Proc. Roy. Soc. Edinburgh Sect. A 115 (1990) no. 3-4, p. 193-230
\bibitem{Tay} M. E. Taylor - Noncommutative harmonic analysis , Math. Surveys and Monographs , vol. 22 , American Mathematical Society, Providence, RI, 1986
\end{thebibliography}
